\section{Finite coorientations eliminate a redundant orbit query}
\label{sec:normal-coorientation}

The normal-surface topology API previously counted the orbits of a supplied
surface, its doubled coordinates, and its boundary separately. Doubling
detects orientability, but its orbit computation can repeat much of the
first query's work. On the 256-tetrahedron layered meridian, the surface and
double each take 261 cycles and record 82,960 local events. The boundary
adds 259 cycles and 518 events. A diagnostic profile identifies transmission,
contraction and support closure in these queries and their independent
replay as the main costs. The profile includes instrumentation overhead;
the timings below come from separate isolated measurements.

The implementation now first tries a sufficient finite coorientation
certificate. When it succeeds, the double has exactly twice the already
verified number of surface components. This removes a complete orbit query
and its local trace. Failure retains the previous query and proof format.
The test is enabled by default in \code{normal\_surface\_topology};
\code{coorientation=False} restores the old schedule. This API analyzes
supplied normal coordinates. The current diagram-recognition stages use
other specialized surface APIs, so these measurements do not assert a
speedup of complete knot recognition.

\subsection{A finite sufficient trivialization}

Let $Q$ be the normal vector used by the orbit queries. If multiplicity
reduction is enabled, the input is $gQ$, where $g$ is the coordinate gcd;
otherwise $Q$ is the full input. Order the global edges as in the normal
interval construction and let $w_e$ be their intersection weights. The
orbit universe is partitioned into nonempty blocks
\[
 B_e=[s_e,s_e+w_e)\cap\mathbb Z,\qquad
 s_e=\sum_{e'<e}w_{e'}.
\]
Each original normal-arc pairing has its domain in one block $B_u$ and
its range in one block $B_v$. Associate to this pairing the bit
$\delta=0$ for an order-preserving map and $\delta=1$ for a reversing map.
Solve the finite graph equations
\[
 \epsilon_u\mathbin\oplus\epsilon_v=\delta.
\]
A parity traversal either supplies one bit per incident block or finds an
inconsistent cycle. The latter is only a failed sufficient test. It does
not prove the surface nonorientable, since a valid coorientation might
vary within a block.

To prove the positive case, write each point of the doubled universe
uniquely as $2i+r$, where $r\in\{0,1\}$. Doubling an interval pairing
$f$ sends it by
\[
 \widetilde f(2i+r)=2f(i)+(r\mathbin\oplus\delta).
\]
This follows directly from the doubled endpoints: an interval $[a,b]$
becomes $[2a,2b+1]$. On block $B_e$, give this point the corrected label
$r\mathbin\oplus\epsilon_e$. The equations imply that every doubled
pairing preserves this label. For each label, projection to $i$ is a
bijection and conjugates every generating pairing to its base pairing.
Thus each label carries exactly the original orbit relation, proving
\[
 C(2Q)=2C(Q).
\]
Blocks with no incident pairing can be given bit zero: their isolated
points have two isolated lifts. The equality also holds for the empty
universe. Since the validated ambient manifold is orientable and this
double is the normal orientation cover, every component of $Q$ is
orientable.

The proof must use the original arc pairings, including reversals on
single-point intervals. Such a reversal is the identity on its base
point but interchanges the two lifted points. Removing its reversal bit
would lose the monodromy and could give an incorrect orientability result.
The graph is therefore reconstructed from the original normal geometry,
independently of simplifications made during the first orbit query.

Coarse-test failure occurs on both one-sided and two-sided surfaces.
The one-tetrahedron M\"obius band fails; twice that band is an orientable
annulus but can still fail, including when its primitive quotient is used.
The one-vertex linking disc also fails on this example. All use the old
double query. When the test succeeds on a quotient $Q$, the existing
multiplicity-lifting identities recover the topology of $gQ$. No new
rule replaces those identities, including the even-multiplicity handling
of one-sided components on the fallback path.

\subsection{An independently checked derived count}

Successful proofs use outer schema \code{normal-surface-topology-v4}.
It retains the source fingerprint, the actual surface and boundary orbit
proofs, boundary homology, claimed topology and coordinate divisor, and
adds an explicit Boolean boundary-classification flag. The double entry
has schema \code{normal-double-coorientation-v1}, the finite parity
certificate, the claimed doubled count and an empty operations list.
The coordinate divisor is explicit even when it equals one.

The checker reconstructs the manifold, admissible coordinates, quotient
and original arc system. It independently verifies the surface count,
reconstructs the finite graph, checks all supplied bits and graph equations,
and requires the double's count to equal twice that verified count. It
does not invoke the parity search or any orbit producer. Missing or
duplicate labels, values other than integer bits, inconsistent equations, extra fields,
nonempty derived traces and incorrect counts are rejected. A simultaneous
bit reversal is a valid alternative certificate. Existing versions one
through three remain accepted.

Boundary classification continues to use its existing component identities
and, when needed, cone queries. Both producer and checker still construct
the doubled arc system, which those optional queries may need. Only its
redundant orbit count is skipped. Failed finite tests produce exactly the
old complete result, including its certificate and statistics. Successful
queries report zero double cycles and a derived-coorientation statistic.
Cycle allowances count actual orbit cycles; replay-operation allowances
count recorded local events. Neither allowance purports to bound all
geometry or graph work. Cancellation callbacks remain active in that work.

For $N$ tetrahedra there are $O(N)$ edge blocks and original arc pairings.
The current helper sorts edge roots, computes binary offsets, and locates
each pairing's four endpoints by binary search. It requires
$O(N\log N)$ integer operations and comparisons, followed by a linear
parity traversal. If the endpoint bit length is bounded by $B$, a safe
elementary bit bound is $O(N\log N(B+\log N))$. Storage is $O(N)$ integers;
the color witness has $O(N\log N)$ bits. No sheet multiplicity is expanded.
These bounds describe the additional test, not the remaining orbit
queries or discovery of a useful normal vector.

\subsection{Validation, costs and scope}

The frozen comparison is revision \code{7d9a6f608}, loading its producer
and independent checker separately. The maintained corpus supplies 1,275
normal vectors on 45 triangulations, including compatible mixtures,
parallel surfaces, empty vectors, one-sided components and relabelled
examples. For each vector the audit tests both multiplicity settings and
both boundary-classification settings: 5,100 configurations. Disabling
the shortcut reproduces all old records exactly. Enabling it derives 956
double counts, corresponding to 239 input vectors, and leaves all 4,144
fallback records identical. Every topology result agrees with the old
result. The audit replays each old proof with both checkers and each new
proof with the new checker, for 15,300 comparisons. Live Regina checks
component counts, orientability, boundary counts and Euler characteristic
on all 1,275 vectors; the frozen component-level records also check the
finer boundary classification.

Focused tests verify the pointwise doubled-map conjugacy, relabellings,
empty inputs, binary multiplicities exceeding 20,000 bits, exact cycle
and replay-event limits, cancellation and proof mutations. Replay is also
run with both orbit search and parity search disabled. All 46 focused
tests and all 1,250 maintained tests pass. The audit retains 84 representative
full proofs, source hashes and deterministic operation counts.
The layered meridian family illustrates the saved work. The table counts
all surface, double and boundary events; the finite sign witness replaces
only the double trace.
\begin{center}
\small
\begin{tabular}{rrrrrr}
\toprule
$N$ & old events & new events & sign bits & old bytes & new bytes\\
\midrule
4 & 118 & 66 & 6 & 8341 & 5038\\
16 & 838 & 438 & 18 & 64708 & 34023\\
64 & 10918 & 5526 & 66 & 907076 & 458720\\
256 & 166438 & 83478 & 258 & 14185011 & 7114254\\
\bottomrule
\end{tabular}
\end{center}


The isolated timing experiment uses five randomized rounds, with old and
new implementations and a same-implementation control for each. All 240
measured calls and 48 warm-up calls finish. Each timed call includes topology
production, independent replay and complete proof serialization and hashing.
Input construction is outside the timer. The old arm uses the old checker;
all arms return the same topology and each arm's proof is deterministic.
The binary-scale case multiplies a 64-tetrahedron meridian by $2^{2048}+1$.
The empty, M\"obius, linking, mixture and genus-two cases also request
boundary classification.
\begin{center}
\small
\begin{tabular}{lrrrrr}
\toprule
Supplied surface & old ms & new ms & old/new & old A/A & new A/A\\
\midrule
Layered 1 & 1.113 & 0.936 & 1.191 & 1.028 & 1.000\\
Layered 16 & 18.532 & 11.317 & 1.638 & 1.009 & 1.000\\
Layered 64 & 227.808 & 121.452 & 1.859 & 0.978 & 1.015\\
Layered 256 & 3490.381 & 1772.576 & 1.967 & 1.002 & 1.000\\
Layered 64, binary scale & 221.309 & 122.368 & 1.810 & 1.017 & 1.027\\
Empty & 0.442 & 0.468 & 0.939 & 0.959 & 1.030\\
M\"obius band & 0.719 & 0.747 & 0.962 & 0.980 & 0.984\\
Doubled M\"obius band & 0.730 & 0.757 & 0.974 & 1.003 & 0.969\\
Vertex-link disc & 1.548 & 1.570 & 0.983 & 1.006 & 1.003\\
Sphere/disc/one-sided mixture & 5.509 & 5.534 & 0.999 & 0.985 & 0.996\\
Sphere plus negative $\chi$ & 5.596 & 5.612 & 0.992 & 0.999 & 0.996\\
Genus-two coherent & 5.398 & 3.995 & 1.323 & 1.020 & 1.001\\
Entire 1,275-vector batch & 5025.707 & 4988.370 & 1.003 & 0.992 & 0.996\\
\bottomrule
\end{tabular}
\end{center}


A factor above one favors the new implementation; factors are medians of
paired ratios rather than ratios of medians. Layered examples gain up to
$1.967\times$, and the genus-two coherent example gains $1.323\times$.
Several small cases regress, including the empty surface even though its
double is derived: the extra graph/proof overhead exceeds the saved query.
Same-implementation control medians range from 0.959 to 1.030, so small
differences should not be interpreted as stable gains.

A second isolated experiment times the entire 1,275-vector corpus as one
batch, with default multiplicity reduction and no boundary classification.
Five rounds give 20 measured batches (25,500 surface calls) and four warm-up
batches (5,100 calls). Each includes all producers, independent replays
and proof serialization/hashing. The aggregate topology digest agrees
across all arms. The new arm derives 239 doubles, reduces recorded events
from 190,570 to 182,925 and serialized proof bytes from 13,330,351 to
12,861,213. Its paired speed factor is 1.003, with old and new control
factors 0.992 and 0.996: effectively neutral at this measurement resolution.
The default therefore trades small miss costs for substantial savings on
some expensive supplied surfaces; it is not presented as a measured
whole-corpus speedup. The explicit off switch preserves the old schedule.


The native research driver is
\path{../fast/normal_orbit_research/coorientation.py}, with \code{audit}
and \code{benchmark} modes, a required \code{--output} path and optional
\code{--rounds}. The separate batch driver is
\path{data/coorientation_corpus_benchmark.py}. Records and validation logs
are \path{data/coorientation-*}. These results justify the finite shortcut
as a way to remove demonstrably redundant work when a small certificate
exists. They do not establish a general speedup on the corpus. The base
orbit search, candidate discovery and authentication back to a knot diagram
remain separate obligations. General quasipolynomial or subexponential
recognition is not proved by this change.
